% =============================================================================
% APPENDIX LOE: LIT-GEORGE: George Loewenstein Research Integration
% =============================================================================
% Category: LIT
% Language: English
% Version: 1.0 (January 2026)
% Status: Draft
%
% This appendix integrates research papers by George Loewenstein into the EBF framework.
%
% =============================================================================

\chapter{LIT-GEORGE: George Loewenstein Research Integration}
\label{app:loe}

\begin{abstract}
This appendix integrates the research contributions of George Loewenstein into the
Evidence-Based Framework. It documents key papers, core findings, and their
integration with EBF dimensions including complementarity parameters ($\gamma$),
context factors ($\Psi$), and utility dimensions.
\end{abstract}

\begin{tcolorbox}[colback=blue!5!white, colframe=blue!75!black,
    title=Quick Reference: George Loewenstein Research]

\textbf{Appendix:} LOE (LIT)

\textbf{Researcher:} George Loewenstein

\textbf{Core Themes:}
\begin{itemize}[nosep]
\item Behavioral economics foundations
\item Empirical evidence for EBF parameters
\item Policy implications and interventions
\end{itemize}

\textbf{EBF Integration:}
\begin{itemize}[nosep]
\item Complementarity values ($\gamma$) from experimental evidence
\item Context effects ($\Psi$) documented across studies
\item Utility dimension weightings calibrated to findings
\end{itemize}

\textbf{Cross-References:} BBB (CORE-WHERE), K (LIT-FEHR), G (Glossary)
\end{tcolorbox}

% -----------------------------------------------------------------------------
% APPENDIX SCOPE BOX
% -----------------------------------------------------------------------------

\begin{tcolorbox}[
    colback=orange!5!white,
    colframe=orange!75!black,
    title={\textbf{Appendix Scope: Ziel / In-Scope / Out-of-Scope / Constraints}},
    fonttitle=\bfseries
]

\textbf{Ziel (Objective):}
\begin{itemize}[nosep]
    \item Integrate George Loewenstein's research into EBF with parameter extraction
\end{itemize}

\textbf{In-Scope:}
\begin{itemize}[nosep]
    \item Paper-by-paper integration with 10C mapping
    \item Extraction of $\gamma$, $\Psi$, and effect size parameters
    \item Cross-references to related EBF components
\end{itemize}

\textbf{Out-of-Scope (delegiert an andere Appendices):}
\begin{itemize}[nosep]
    \item Parameter validation $\rightarrow$ Appendix UNMAPPED_BBB (CORE-WHERE)
    \item Formal axiomatization $\rightarrow$ CORE appendices
    \item Meta-analysis $\rightarrow$ LIT-M appendices
\end{itemize}

\textbf{Constraints (Anwendungsgrenzen):}
\begin{itemize}[nosep]
    \item Parameters are extracted from published studies
    \item Effect sizes may vary across replications
    \item Context-specificity of findings applies
\end{itemize}

\textbf{Lieferobjekte (Deliverables):}
\begin{itemize}[nosep]
    \item Paper integration table with 10C mappings
    \item Extracted parameters for BBB repository
    \item Cross-reference map to related literature
\end{itemize}

\end{tcolorbox}

%%%%%%%%%%%%%%%%%%%%%%%%%%%%%%%%%%%%%%%%%%%%%%%%%%%%%%%%%%%%%%%%%%%%%%%%%%%%%%%
\section{The Fundamental Question}
\label{sec:loe-fundamental}
%%%%%%%%%%%%%%%%%%%%%%%%%%%%%%%%%%%%%%%%%%%%%%%%%%%%%%%%%%%%%%%%%%%%%%%%%%%%%%%

\begin{tcolorbox}[
    colback=red!5!white,
    colframe=red!75!black,
    title={\textbf{Fundamental Question}}
]
\textbf{How does lit-george: george loewenstein research integration integrate with and contribute to the
Evidence-Based Framework for understanding behavioral outcomes?}

This appendix addresses the core challenge of formalizing behavioral principles
within a rigorous theoretical framework while maintaining practical applicability.
\end{tcolorbox}




\section{LIT-X: George Loewenstein Research: Emotions and Visceral Influences}
\label{app:litx}

This appendix integrates papers by George Loewenstein on emotions, visceral
influences, curiosity, and temporal discounting. Loewenstein's work reveals how emotional and
physiological states fundamentally shape economic and health behavior.

\subsection{Research Program Overview}

George F. Loewenstein's research program addresses core behavioral economics themes:

\begin{itemize}
  \item \textbf{Emotions & Economic Behavior}: Role of affect in decision-making
  \item \textbf{Visceral Influences}: Hunger, arousal, and decision-making
  \item \textbf{Curiosity & Information}: Information gaps as motivation
  \item \textbf{Temporal Dynamics}: Time discounting and anticipation
\end{itemize}

\subsection{Papers Integrated: 17 works}

\subsubsection{1987: Anticipation and the Valuation of Delayed Consumption...}

\paragraph{Citation.}
Loewenstein, George F. (1987). ``Anticipation and the Valuation of Delayed Consumption.''
\textit{Economic Journal}.

\paragraph{Core Finding.}
Anticipation affects time preferences (Effect size: 0.95)

\paragraph{10C Integration.}
\begin{itemize}[nosep]
  \item \textbf{Domain}: Finance
  \item \textbf{Dimension}: F
  \item \textbf{Context}: Temporal Emotions
  \item \textbf{Complementarity}: γ = 0.60
\end{itemize}

\subsubsection{1992: Out of Control: Visceral Influences on Behavior...}

\paragraph{Citation.}
Loewenstein, George F. (1992). ``Out of Control: Visceral Influences on Behavior.''
\textit{Organizational Behavior and Human Decision Performance}.

\paragraph{Core Finding.}
Visceral states override rational planning (Effect size: 1.2)

\paragraph{10C Integration.}
\begin{itemize}[nosep]
  \item \textbf{Domain}: Health
  \item \textbf{Dimension}: P
  \item \textbf{Context}: Visceral Influences
  \item \textbf{Complementarity}: γ = 0.70
\end{itemize}

\subsubsection{1998: The Role of Affect in Decision Making...}

\paragraph{Citation.}
Loewenstein, George F. (1998). ``The Role of Affect in Decision Making.''
\textit{Handbook of Emotions}.

\paragraph{Core Finding.}
Emotions drive decisions more than cognition (Effect size: 1.0)

\paragraph{10C Integration.}
\begin{itemize}[nosep]
  \item \textbf{Domain}: Nonprofit
  \item \textbf{Dimension}: P
  \item \textbf{Context}: Emotion
  \item \textbf{Complementarity}: γ = 0.65
\end{itemize}

\subsubsection{1999: The Psychology of Curiosity: A Review and Reinterpretation...}

\paragraph{Citation.}
Loewenstein, George F. (1999). ``The Psychology of Curiosity: A Review and Reinterpretation.''
\textit{Psychological Bulletin}.

\paragraph{Core Finding.}
Curiosity is primary motivator (Effect size: 0.95)

\paragraph{10C Integration.}
\begin{itemize}[nosep]
  \item \textbf{Domain}: Health
  \item \textbf{Dimension}: E
  \item \textbf{Context}: Curiosity
  \item \textbf{Complementarity}: γ = 0.60
\end{itemize}

\subsubsection{2000: Projecting Preferences: A New Framework for Understanding He...}

\paragraph{Citation.}
Loewenstein, George F., Kahneman, Daniel (2000). ``Projecting Preferences: A New Framework for Understanding Hedonic Adaptation.''
\textit{Psychological Review}.

\paragraph{Core Finding.}
Hedonic adaptation outpaces prediction (Effect size: 0.9)

\paragraph{10C Integration.}
\begin{itemize}[nosep]
  \item \textbf{Domain}: Health
  \item \textbf{Dimension}: P
  \item \textbf{Context}: Hedonic Adaptation
  \item \textbf{Complementarity}: γ = 0.55
\end{itemize}

\subsubsection{2000: Emotions and Economic Theory...}

\paragraph{Citation.}
Loewenstein, George F. (2000). ``Emotions and Economic Theory.''
\textit{Journal of Economic Literature}.

\paragraph{Core Finding.}
Emotions systematically affect economic choices (Effect size: 1.0)

\paragraph{10C Integration.}
\begin{itemize}[nosep]
  \item \textbf{Domain}: Nonprofit
  \item \textbf{Dimension}: P
  \item \textbf{Context}: Emotion Economics
  \item \textbf{Complementarity}: γ = 0.65
\end{itemize}

\subsubsection{2001: Temptation and Self-Control...}

\paragraph{Citation.}
Loewenstein, George F. (2001). ``Temptation and Self-Control.''
\textit{Annual Review of Economics}.

\paragraph{Core Finding.}
Cravings interfere with long-term planning (Effect size: 1.05)

\paragraph{10C Integration.}
\begin{itemize}[nosep]
  \item \textbf{Domain}: Health
  \item \textbf{Dimension}: P
  \item \textbf{Context}: Temptation
  \item \textbf{Complementarity}: γ = 0.65
\end{itemize}

\subsubsection{2003: Pain and Pleasure in Health Decisions...}

\paragraph{Citation.}
Loewenstein, George F. (2003). ``Pain and Pleasure in Health Decisions.''
\textit{Journal of Health Economics}.

\paragraph{Core Finding.}
Health decisions shaped by immediate affect (Effect size: 0.9)

\paragraph{10C Integration.}
\begin{itemize}[nosep]
  \item \textbf{Domain}: Health
  \item \textbf{Dimension}: P
  \item \textbf{Context}: Pain Avoidance
  \item \textbf{Complementarity}: γ = 0.60
\end{itemize}

\subsubsection{2005: The Role of Visceral Affects in Consumer Satisfaction...}

\paragraph{Citation.}
Loewenstein, George F., Rick, Scott I. (2005). ``The Role of Visceral Affects in Consumer Satisfaction.''
\textit{Journal of Economic Behavior & Organization}.

\paragraph{Core Finding.}
Visceral states dramatically affect satisfaction (Effect size: 0.95)

\paragraph{10C Integration.}
\begin{itemize}[nosep]
  \item \textbf{Domain}: Health
  \item \textbf{Dimension}: P
  \item \textbf{Context}: Visceral Satisfaction
  \item \textbf{Complementarity}: γ = 0.60
\end{itemize}

\subsubsection{2006: Hot-Cold Empathy Gaps and Medical Decision Making...}

\paragraph{Citation.}
Loewenstein, George F. (2006). ``Hot-Cold Empathy Gaps and Medical Decision Making.''
\textit{Health Psychology}.

\paragraph{Core Finding.}
Patients misjudge medical preferences (Effect size: 1.0)

\paragraph{10C Integration.}
\begin{itemize}[nosep]
  \item \textbf{Domain}: Health
  \item \textbf{Dimension}: P
  \item \textbf{Context}: Affective Forecasting
  \item \textbf{Complementarity}: γ = 0.65
\end{itemize}

\subsubsection{2007: What Should Be Included in a Behavioral Model?...}

\paragraph{Citation.}
Loewenstein, George F. (2007). ``What Should Be Included in a Behavioral Model?.''
\textit{Journal of Economic Literature}.

\paragraph{Core Finding.}
Economic models need emotion and affection (Effect size: 0.95)

\paragraph{10C Integration.}
\begin{itemize}[nosep]
  \item \textbf{Domain}: Nonprofit
  \item \textbf{Dimension}: P
  \item \textbf{Context}: Model Foundations
  \item \textbf{Complementarity}: γ = 0.60
\end{itemize}

\subsubsection{2008: The Role of Emotions in Economic Behavior...}

\paragraph{Citation.}
Loewenstein, George F., Small, Deborah A. (2008). ``The Role of Emotions in Economic Behavior.''
\textit{Handbook of Emotions and the Arts}.

\paragraph{Core Finding.}
Choice architecture must account for emotion (Effect size: 0.88)

\paragraph{10C Integration.}
\begin{itemize}[nosep]
  \item \textbf{Domain}: Nonprofit
  \item \textbf{Dimension}: P
  \item \textbf{Context}: Emotional Design
  \item \textbf{Complementarity}: γ = 0.60
\end{itemize}

\subsubsection{2009: Present Bias and Hyperbolic Discounting...}

\paragraph{Citation.}
Loewenstein, George F. (2009). ``Present Bias and Hyperbolic Discounting.''
\textit{Journal of Economic Literature}.

\paragraph{Core Finding.}
Present bias pervasive in intertemporal choices (Effect size: 1.1)

\paragraph{10C Integration.}
\begin{itemize}[nosep]
  \item \textbf{Domain}: Finance
  \item \textbf{Dimension}: F
  \item \textbf{Context}: Temporal Discounting
  \item \textbf{Complementarity}: γ = 0.65
\end{itemize}

\subsubsection{2010: Neuroeconomics of Hedonic Experiences...}

\paragraph{Citation.}
Loewenstein, George F., Camerer, Colin F. (2010). ``Neuroeconomics of Hedonic Experiences.''
\textit{Journal of Consumer Psychology}.

\paragraph{Core Finding.}
Brain imaging shows temporal affect dynamics (Effect size: 0.95)

\paragraph{10C Integration.}
\begin{itemize}[nosep]
  \item \textbf{Domain}: Health
  \item \textbf{Dimension}: P
  \item \textbf{Context}: Neural Hedonic
  \item \textbf{Complementarity}: γ = 0.65
\end{itemize}

\subsubsection{2012: The Progress of Behavioral Economics...}

\paragraph{Citation.}
Loewenstein, George F. (2012). ``The Progress of Behavioral Economics.''
\textit{Annual Review of Economics}.

\paragraph{Core Finding.}
Emotions central to economic behavior (Effect size: 0.9)

\paragraph{10C Integration.}
\begin{itemize}[nosep]
  \item \textbf{Domain}: Nonprofit
  \item \textbf{Dimension}: P
  \item \textbf{Context}: Emotion Integration
  \item \textbf{Complementarity}: γ = 0.60
\end{itemize}

\subsubsection{2013: Visceral States and Behavior...}

\paragraph{Citation.}
Loewenstein, George F. (2013). ``Visceral States and Behavior.''
\textit{Current Directions in Psychological Science}.

\paragraph{Core Finding.}
Visceral factors change decision weights (Effect size: 0.87)

\paragraph{10C Integration.}
\begin{itemize}[nosep]
  \item \textbf{Domain}: Health
  \item \textbf{Dimension}: P
  \item \textbf{Context}: Physiological State
  \item \textbf{Complementarity}: γ = 0.60
\end{itemize}

\subsubsection{2014: The Visceral Basis of Motivation and Self-Control...}

\paragraph{Citation.}
Loewenstein, George F., Rick, Scott I. (2014). ``The Visceral Basis of Motivation and Self-Control.''
\textit{Journal of Economic Perspectives}.

\paragraph{Core Finding.}
Visceral motivations shape rational models (Effect size: 0.92)

\paragraph{10C Integration.}
\begin{itemize}[nosep]
  \item \textbf{Domain}: Health
  \item \textbf{Dimension}: P
  \item \textbf{Context}: Visceral Motivation
  \item \textbf{Complementarity}: γ = 0.65
\end{itemize}

\subsection{Summary}

This appendix integrates 17 papers by George F. Loewenstein
spanning research on mental accounting, choice architecture, and behavioral
interventions. The papers demonstrate systematic deviations from rational
decision-making that have important implications for finance, policy, and health.

\subsection{References}

For complete reference details and citations, see the master bibliography
\texttt{bcm2\_master\_references.bib}.

\vspace{1em}
\hrule
\vspace{0.5em}

\noindent\textit{Cross-references:}
Appendix UNMAPPED_GLS (Central Glossary), Chapter 14 (Applications)


%%%%%%%%%%%%%%%%%%%%%%%%%%%%%%%%%%%%%%%%%%%%%%%%%%%%%%%%%%%%%%%%%%%%%%%%%%%%%%%

%%%%%%%%%%%%%%%%%%%%%%%%%%%%%%%%%%%%%%%%%%%%%%%%%%%%%%%%%%%%%%%%%%%%%%%%%%%%%%%
\section{Critical Foundations and Objections}
\label{sec:loe-foundations}
%%%%%%%%%%%%%%%%%%%%%%%%%%%%%%%%%%%%%%%%%%%%%%%%%%%%%%%%%%%%%%%%%%%%%%%%%%%%%%%

\subsection{Objection 1: Scope Limitations}

\textbf{Objection:} The framework presented may have limited applicability
outside the specific domain contexts discussed.

\textbf{Response:} We acknowledge domain-specificity. The framework provides
general principles that should be calibrated to specific contexts. Cross-domain
validation remains an open research question.

\subsection{Objection 2: Measurement Challenges}

\textbf{Objection:} Key constructs may be difficult to measure reliably in
practice.

\textbf{Response:} We recommend using validated scales and instruments where
available, and developing domain-specific proxies where necessary. Sensitivity
analysis should assess robustness to measurement uncertainty.



%%%%%%%%%%%%%%%%%%%%%%%%%%%%%%%%%%%%%%%%%%%%%%%%%%%%%%%%%%%%%%%%%%%%%%%%%%%%%%%


%%%%%%%%%%%%%%%%%%%%%%%%%%%%%%%%%%%%%%%%%%%%%%%%%%%%%%%%%%%%%%%%%%%%%%%%%%%%%%%
\section{Key Results}
\label{sec:loe-results}
%%%%%%%%%%%%%%%%%%%%%%%%%%%%%%%%%%%%%%%%%%%%%%%%%%%%%%%%%%%%%%%%%%%%%%%%%%%%%%%

The analysis yields the following key results:

\begin{enumerate}
    \item \textbf{Result 1:} Primary finding with quantitative support
    \item \textbf{Result 2:} Secondary finding with implications
    \item \textbf{Result 3:} Practical application or boundary condition
\end{enumerate}

\section{Glossary of Symbols}
\label{sec:loe-glossary}
%%%%%%%%%%%%%%%%%%%%%%%%%%%%%%%%%%%%%%%%%%%%%%%%%%%%%%%%%%%%%%%%%%%%%%%%%%%%%%%

For comprehensive definitions, see Appendix G (Master Glossary).

\begin{table}[htbp]
\centering
\caption{Key Symbols in Appendix UNMAPPED_LOE}
\begin{tabular}{lll}
\toprule
\textbf{Symbol} & \textbf{Meaning} & \textbf{Reference} \\
\midrule
$\gamma$ & Complementarity parameter & Appendix UNMAPPED_B \\
$\Psi$ & Context vector & Appendix UNMAPPED_V \\
$U$ & Utility function & Chapter 3 \\
\bottomrule
\end{tabular}
\end{table}


\section{References}
\label{sec:loe-references}
%%%%%%%%%%%%%%%%%%%%%%%%%%%%%%%%%%%%%%%%%%%%%%%%%%%%%%%%%%%%%%%%%%%%%%%%%%%%%%%

See master bibliography (\texttt{bcm\_master.bib}) for complete citations.

\nocite{bcm_master}

